\subsubsection{Show Applicant Profile}
This particular use case will enable the Registrar Office to see the entire profile of any applicant of choice. Once the search for the applicant has been carried out, the user will then be able to click on the name of the applicant or the “View Profile” button to see all information about him/her in one particular page known as the profile. This use case is very useful especially when processing an applicant’s application.Moreover, the profile layout is designed in a manner that organizes the information in categories such as personal information, academic information, and documents, making it easier to read. The system enables the Registrar’s Office to evaluate all the information provided without the need to navigate from one webpage to another. The system provides a proper arrangement of all documents submitted in support of the application, making it easier to verify all documents. Additionally, all changes made to the information provided are updated instantly on the applicant’s profile page.
\begin{table}[H]
    \centering
    \small
    \renewcommand{\arraystretch}{1.4}
    \caption{Use Case Table: Show Applicant Profile (UC-5.2)} % ক্যাপশন উপরে আনা হয়েছে
    \label{tab:show-applicant-profile-registrar}
    \begin{tabularx}{\textwidth}{|l|X|}
        \hline
        \rowcolor[gray]{0.9} {Characteristics} & {Description} \\ \hline
        Use Case ID & UC-5.2 \\ \hline
        Use Case Name & Show Applicant Profile \\ \hline
        Primary Actors & Registrar Office \\ \hline
        Goal & To see the entire personal and academic profile of a particular applicant. \\ \hline
        Precondition & The user is required to look up a certain applicant in the database. \\ \hline
        Scenario & 1. Look up a certain applicant. 2. Click on either the “View Profile” link or the name of the applicant. 3. The system presents the full profile page. \\ \hline
        Exception & No authorization or no matching record. \\ \hline
        Frequency of Use & High frequency use in the data and admission process. \\ \hline
    \end{tabularx}
\end{table}
\FloatBarrier